
\subsubsection*{2.4.B}
\phantomsection
\label{sga3:VIIB:2.4.B}

Suppose, for simplicity, that \(k\) is Artinian. The preceding
discussion can then be expressed more simply in terms of the
covariant bialgebras of \(G\), \(Q\), and \(N\). Indeed, since
\[
  G=N\times Q
\]
as formal \(k\)-varieties, one has
\[
  \mathbf H(G)
  =
  \mathbf H(N)\otimes_k\mathbf H(Q)
\]
as \(k\)-coalgebras. Moreover, since multiplication in \(G\) is
given by
\[
  (x,q)\cdot(x',q')
  =
  \bigl(x\alpha(q,x'),qq'\bigr),
  \qquad
  \alpha(q,x')=qx'q^{-1},
\]
multiplication in \(\mathbf H(G)\) is given, for
\(x\in\mathbf H(N)\) and \(q\in\mathbf H(Q)\), by
\[
  (x\otimes q)\cdot(x'\otimes q')
  =
  x\phi(q,x')\otimes qq',
\]
where
\[
  \phi:
  \mathbf H(Q)\otimes_k\mathbf H(N)
  \longrightarrow
  \mathbf H(N)
\]
is the morphism of \(k\)-coalgebras induced by \(\alpha\).

The morphism \(\alpha\) is the following composite. Here
\(\delta_G\) and \(m_G\) denote, respectively, the diagonal and
multiplication morphisms of \(G\), \(c_Q\) is inversion in \(Q\),
and
\[
  v(q\otimes q'\otimes x)=q\otimes x\otimes q'.
\]
\begin{center}
  \phantomsection
  \label{sga3:VIIB:2.4.B:conjugation-composite}
  \(
Q\times N
\xrightarrow{\ v\circ(\delta_G\times\operatorname{id})\ }
Q\times N\times Q
\xrightarrow{\ \operatorname{id}\times\operatorname{id}\times c_Q\ }
Q\times N\times Q
\xrightarrow{\ m_G\ }G,
\)

\end{center}
Writing \(c_Q\) also for the antipode of \(\mathbf H(Q)\), we obtain
\[
  \phi(q\otimes x')
  =
  \sum_i q_i x'c_Q(q_i')
  \quad\text{if}\quad
  \Delta_{\mathbf H(Q)}(q)
  =
  \sum_i q_i\otimes q_i'.
  \tag{\(\star\)}\label{sga3:VIIB:2.4.B:equation:covariant-antipode-star}
\]

In particular, let \(M\) be an abstract group and suppose that
\(\mathbf H(Q)\) is the group algebra \(kM\), so that
\[
  Q=\operatorname{Spf}^{*}kM
\]
is the constant formal \(k\)-group \(M_k\). Then for every
\(\gamma\in M\) and \(x'\in\mathbf H(N)\),
\[
  \phi(\gamma\otimes x')
  =
  \gamma x'\gamma^{-1}.
\]
This defines an action of \(M\) on \(\mathbf H(N)\) by Hopf-algebra
automorphisms.
